\documentclass[11pt,a4paper]{article}
\usepackage{amsfonts,amsmath,amssymb}
\usepackage{amsthm}
\usepackage[swedish, english]{babel}
\usepackage[T1]{fontenc}
\usepackage[utf8]{inputenc}
\usepackage{comment}
\usepackage[counterclockwise]{rotating}
\usepackage{float}
\usepackage{tikz}
\usepackage{caption}
\usepackage{graphicx}
\usepackage{geometry}
\usepackage{curves}
\usepackage{booktabs, dcolumn}
\usepackage{url}
\usepackage{multirow}
\usepackage{threeparttable}
\usepackage{longtable}
%\usepackage{pstricks}
\usepackage{lastpage}
\usepackage{appendix}
\usepackage{lscape}
\usepackage{color,soul}
\usepackage[dvipsnames]{xcolor}
\usepackage[round,authoryear]{natbib}
\oddsidemargin 0in
\evensidemargin\oddsidemargin
\setlength{\topmargin}{-0.5in}
\setlength{\textheight}{9.5in}
\setlength{\textwidth}{6.3in}
\usepackage{setspace}
\onehalfspacing
\usepackage[parfill]{parskip}
\usepackage[pdftex,breaklinks,colorlinks,
citecolor=blue,
urlcolor=blue,
anchorcolor=blue,
allcolors=blue,
linkcolor=black]{hyperref}
\usepackage{etoolbox}
\makeatletter
\patchcmd{\@figurepage}{\vspace{20pt}}{\clearpage}{}{}
\patchcmd{\@tablepage}{\bigskip}{\clearpage}{}{}
\makeatother


%%%%%%%%%%%%%%%%%%THOMAS RELATED THINGS%%%%%%%%%%%%%%%%%
\setcounter{MaxMatrixCols}{30}
%TCIDATA{OutputFilter=latex2.dll}
%TCIDATA{Version=5.50.0.2960}
%TCIDATA{Codepage=1252}
%TCIDATA{CSTFile=LaTeX article (bright).cst}
%TCIDATA{Created=Thursday, November 06, 2008 11:22:05}
%TCIDATA{LastRevised=Tuesday, March 08, 2016 16:18:15}
%TCIDATA{<META NAME="GraphicsSave" CONTENT="32">}
%TCIDATA{<META NAME="SaveForMode" CONTENT="1">}
%TCIDATA{BibliographyScheme=Manual}
%TCIDATA{<META NAME="DocumentShell" CONTENT="Articles\SW\Standard LaTeX Article (Harvard)">}
%TCIDATA{Language=American English}
%BeginMSIPreambleData
\providecommand{\U}[1]{\protect\rule{.1in}{.1in}}
%EndMSIPreambleData
%\usepackage[pdftex,breaklinks,colorlinks,
%citecolor=blue,
%urlcolor=blue,
%linkcolor=black]{hyperref}
\newtheorem{theorem}{Theorem}
\newtheorem{acknowledgement}[theorem]{Acknowledgement}
\newtheorem{algorithm}[theorem]{Algorithm}
\newtheorem{axiom}[theorem]{Axiom}
\newtheorem{case}[theorem]{Case}
\newtheorem{claim}[theorem]{Claim}
\newtheorem{conclusion}[theorem]{Conclusion}
\newtheorem{condition}[theorem]{Condition}
\newtheorem{conjecture}[theorem]{Conjecture}
\newtheorem{corollary}[theorem]{Corollary}
\newtheorem{criterion}[theorem]{Criterion}
\newtheorem{definition}[theorem]{Definition}
\newtheorem{example}[theorem]{Example}
\newtheorem{exercise}[theorem]{Exercise}
\newtheorem{lemma}{Lemma}
\newtheorem{notation}[theorem]{Notation}
\newtheorem{problem}[theorem]{Problem}
\newtheorem{proposition}{Proposition}
\newtheorem{assumption}{Assumption}
\newtheorem{remark}[theorem]{Remark}
\newtheorem{solution}[theorem]{Solution}
\newtheorem{summary}[theorem]{Summary}
\usepackage{lineno}
%\newenvironment{proof}[1][Proof]{\noindent\textbf{#1.} }{\ \rule{0.5em}{0.5em}}
\geometry{left=2.5cm,right=2.5cm,top=3cm,bottom=3cm}


%%%%%%%%%%%%%%%%ESTOUT RELATED THINGS%%%%%%%%%%%%%%%%%%%
\newcommand{\sym}[1]{\rlap{#1}}% Thanks to David Carlisle

\let\estinput=\input% define a new input command so that we can still flatten the document

\newcommand{\estwide}[3]{
		\vspace{.75ex}{
			\begin{tabular*}
			{\textwidth}{@{\hskip\tabcolsep\extracolsep\fill}l*{#2}{#3}}
			\toprule
			\estinput{#1}
			\bottomrule
			\addlinespace[.75ex]
			\end{tabular*}
			}
		}	

\newcommand{\estauto}[3]{
		\vspace{.75ex}{
			\begin{tabular}{l*{#2}{#3}}
			\toprule
			\estinput{#1}
			\bottomrule
			\addlinespace[.75ex]
			\end{tabular}
			}
		}

% Allow line breaks with \\ in specialcells
	\newcommand{\specialcell}[2][c]{%
	\begin{tabular}[#1]{@{}c@{}}#2\end{tabular}}

% *****************************************************************
% Custom subcaptions
% *****************************************************************
% Note/Source/Text after Tables
\newcommand{\figtext}[1]{
	\vspace{-0.5ex}
	\captionsetup{justification=justified,font=footnotesize}
	\caption*{\hspace{6pt}\hangindent=1.5em #1}
	}
\newcommand{\fignote}[1]{\figtext{{}~#1}}

\newcommand{\figsource}[1]{\figtext{\emph{Source:~}~#1}}

% Add significance note with \starnote
\newcommand{\starnote}{\figtext{* p < 0.1, ** p < 0.05, *** p < 0.01. Standard errors in parentheses.}}

% *****************************************************************
% siunitx
% *****************************************************************
\usepackage{siunitx} % centering in tables
	\sisetup{
		detect-mode,
		tight-spacing		= true,
		group-digits		= false ,
		input-signs		= ,
		input-symbols		= ( ) [ ] - + *,
		input-open-uncertainty	= ,
		input-close-uncertainty	= ,
		table-align-text-post	= false
        }

%%%%%%%%%%%%%%%%%%%%%%%%%%%%%%%%%%%%%%%%%%%%%%%%%%%%
%\graphicspath{ {"/files/graphs/"} 

%%%%%%%%%%WHEN PREPARING FOR SUBMISSION: 
%SET FIGCAPS ON, BUT SET FIGCAPS OFF WHEN APPENDICES BEGIN 


\makeatletter
\def\input@path{{graphs/}}
\makeatother


\graphicspath{{graphs/}}



